\section{Magnetization dynamic of a $\mathrm{CoFe/Co_2MnSi}$ magnetic bilayer structure\href{https://doi.org/10.1063/5.0128519}{[\textit{J. Appl. Phys., 133 093902 (2023)}]}}


\subsection*{\underline{Abstract}}
Magnetization dynamics of a thin epitaxial $\mathrm{Co_2MnSi(CMS)}$ in magnetic $\mathrm{CoFe/CMS}$ bilayer structured sputtered on MgO substrate is studied. Fourfold symmetry with respect to the azimuthal applied field angle was obtained as the frequency response of the magnetization precession of $\mathrm{CoFe/CMS}$. Moreover, the effective Gilbert damping $\alpha_{\mathrm{eff}}$ exhibits inhomogeneous broadening at lower applied magnetic fields. And at large fields, the azimuthal angle dependence disappears, and the intrinsic damping is observed. 

\subsection*{\underline{I. Introduction}}
\begin{enumerate}
    \item Due to the epitaxial growth of magnetic multi-layers, precise control of electronic and magnetic properties for applications such as magnetic random access memory (MRAM), sensors, and spin transistors is possible.
    \item These devices rely mostly on the multi-layer magnetic tunnel junctions (MTJs), in which high spin polarization is essential to maximize the TMR, and end up optimizing the efficiency of the whole system.
    \item Heusler alloys are coming up as very promising candidates because of the following reasons;
    \begin{itemize}
        \item Half-metallic nature,
        \item High Curie temperature, and 
        \item particularly $\mathrm{CMS}$ being one of the most promising as it can crystallize in the ordered $L\mathrm{2_1}$ crystal structure, as theoretically it has been shown to possess zero spin down DOS at the Fermi energy.
        \item Excellent static behavior with a TMR ratio of $570\%$ at low temperature and $753\%$ in $\mathrm{CMS/MgO/CoFe}$ structure.
        \item In MTJ applications, Heusler alloys are also promising materials for THz emitters [\textcolor{red}{ref. 28,29}].
    \end{itemize}
    \item Understanding the dynamic properties provides time-dependent device behavior and can reveal important material parameters, such as the Gilbert damping parameter $\mathrm{\alpha_{eff}}$.
    \item And the lowering of the damping in the case of $\mathrm{Co_2FeAl(CFA)}$ and  $\mathrm{Co_2MnSi(CMS)}$ leads to the suppression of the spin-flip processes.
    \item The measurement of the precession frequency of the CMS film at low applied magnetic fields (below 1kOe) and as a function of the azimuthal angle $\phi$ between the applied field and the [110] crystalline direction of the film.
    \item A fourfold symmetric dependence of both the precession frequency and damping as a function of $\phi$ was observed, caused by the magnetocrystalline anisotropy present in this cubic structure.
    \item Using various magnetic fields applied at various polar and azimuthal angles, they have extracted the magnetization dynamics all the way up to the high-field limit, where inhomogeneous broadening mechanisms disappear, and intrinsic Gilbert damping of the material can be determined.
    \item \hl{Additionally, it has been found that the angular dependence of the damping parameter disappears in the high magnetic field}.
\end{enumerate}

\subsection*{\underline{II. Sample and Experimental Procedure}}
\begin{enumerate}
    \item Two different samples, $\mathrm{Co_{50}Fe_{50}}$ film was prepared on MgO[001] substrate, and the second sample consisted on a CoFe(1.7 nm)/CMS (7.0 nm) bilayer film.
    \item The layer thickness of the CoFe/CMS bilayer film was determined from X-ray reflectivity, and 1.7 nm CoFe and 7.0 nm CMS were recorded with sharp interfaces of \hl{0.3 nm roughness}.
    \item The bulk CMS compound has a \hl{lattice constant of 0.5654 nm} [\textcolor{red}{ref. 39}].
    \item Magnetization dynamics were analyzed using a dual-color time-resolved magneto-optic Kerr effect (TR-MOKE) setup in polar Kerr geometry.
    \item As the pump beam is spectrally filtered after the reflection, the collected signals provide information on the change in polarization due to the magnetization dynamics.
\end{enumerate}
\subsection*{\underline{III. Results and Discussion}}
\begin{enumerate}
    \item The magnetization dynamics of the epitaxial CoFe were first analyzed, with an external magnetic field of $\mathrm{H_{app}=5}$ kOe applied at multiple polar angles $(\theta,$ angle with respect to the z-axis$)$ and various azimuthal angles $(\phi,$ angle with respect to the x-axis in the $x-y$ plane$)$.
    \item The fourfold symmetry in the BCC structure was probed particularly at polar angles of $\theta=30^\circ,45^\circ,$ and $55^\circ$, and the azimuthal angle was scanned in each case between $0^\circ \leq \phi\leq 90^\circ$ to probe the full fourfold symmetry in the BCC structure.
    \item The dependence of the magnetization precession as a function of azimuthal angles is obtained for three different polar angles.
    \item The effective damping parameter was highest when the in-plane component of the magnetic field was applied along the hard axis $(\phi=0^\circ, 90^\circ)$ and lowest when the in-plane component of the magnetic field was applied along the easy axis  $(\phi=45^\circ)$. [\textcolor{blue}{Important line}]
    \item In the case of the single-layer CoFe film, the frequency response was symmetrically centered around $\phi=45^\circ$.
    \item The azimuthal scan now shows a frequency maximum when the in-plane component of the magnetic field was applied along the $\langle110\rangle$ direction and a minimum along the $\langle100\rangle$.
    \item The field-dependent effective damping parameter was extracted by fitting the data using $\theta_k=A \sin(\omega t+\varphi)e^{-t/\tau}$, $\omega=2\pi f$
    \item In this CMS film, two-magnon scattering (TMS), e.g., due to strain-induced anisotropy inhomogeneities, is a dominant source of extrinsic damping.
    \item At low fields, the apparent difference in $\mathrm{\alpha_{eff}}$ as the azimuthal angle is varied, while the high-field limit appears uniform for all orientations.
    \item The clear dependence of $\mathrm{\alpha_{eff}}$ on $\phi$ was observed with maximum damping at $\phi=45^\circ$, in agreement with the observations of Liu $et$ al. where the inhomogeneous broadening explains the magnetocrystalline anisotropy.
    \item At high applied fields, however, this dependence disappears, and the intrinsic Gilbert damping of $\sim 0.02$ is observed for all parameter combinations.
    \item We note that the relatively $\mathrm{\alpha_{eff}}$ may be due to the mechanical strains present in both the very thin CoFe and the CMS films.
    \item And these strains create defects in the crystal which cause additional spin wave scattering, leading to a larger effective damping parameter $\mathrm{\alpha_{eff}}$.
\end{enumerate}
